% Versão provisória (projeto, sem resultados). Reescrever após o Cap. 5:
% contexto · problema · método (B1 + B2) · resultados · conclusão.
% 150–500 palavras (NBR 6028).
A integração de modelos de linguagem de grande escala a aplicações \textit{web}
costuma ser feita por chamadas diretas ao kit de desenvolvimento do provedor,
espalhadas pelo código de negócio, o que acopla a aplicação a um fornecedor
específico. Uma alternativa é interpor uma camada arquitetural intermediária
--- aqui denominada \textit{harness} --- que centraliza a comunicação com os
provedores, a execução de ferramentas, o estado, o tratamento de falhas e o
registro de eventos. A literatura propõe arquiteturas desse tipo e discute seus
compromissos de forma qualitativa, mas não os mede em estudo controlado. Este
trabalho avalia experimentalmente o \textit{trade-off} entre desacoplamento e
eficiência energética introduzido por essa camada, comparando duas versões
funcionalmente equivalentes de uma mesma aplicação \textit{web}: uma integrada
ao provedor por meio do \textit{harness} e outra integrada diretamente. O custo
de substituição do provedor é quantificado por métricas comportamentais e
estruturais do código, e o custo energético, por telemetria de \textit{hardware}
em requisições simples e agênticas a um modelo aberto executado localmente.
Espera-se caracterizar em que condições o desacoplamento obtido compensa o
custo energético da camada, oferecendo evidência quantitativa para uma decisão
arquitetural hoje tomada com base em princípios de projeto ou em alegações de
fornecedores.

% Separe as palavras-chave por ponto
\palavraschave{Arquitetura de software. Integração de LLM. Substituibilidade. Eficiência energética. Green AI.}
